%قبل از شروع مسئله ترجیحاً چیزی ننویسید. اما اگر خواستید بنویسید، در قالب لیستی (فقرات و شمارش) بنویسید.
%برای اضافه کردن نمره، بعدی از شروع‌مسئله، [x] را اضافه کنید و به جای x  نمره مورد نظر را قرار دهید. برای زیرمسئله‌ها هم به همین صورت است.
%\شروع‌مسئله[10]
\شروع‌مسئله[16]
% صورت اصلی مسئله را در این‌جا بنویسید.
فرض کنید $n$تا رشته داریم 
و می‌خواهیم رشته‌ها را به ترتیب الفبایی مرتب‌کنیم.
رشته $s$
در ترتیب الفبایی از 
$p$
کوچکتر است اگر 
$s$
یک پیشوند از 
$p$
باشد یا اینکه اگر 
$i$
کوچکترین مقداری باشد که
$s_i \neq p_i$
در این‌صورت
$s_i < p_i$
باشد.

برای مقایسه دو تا رشته باهم می‌توانیم روی تک تک کاراکترها از اول پیمایش کنیم و هر گاه به اولین جایگاهی رسیدیم که کاراکترهای متناظر نابرابرند آن کاراکتر را مقایسه کنیم.

حالا فرض کنید الگوریتم 
\LRE{merge sort}
را روی این تعداد کاراکتر اجرا می‌کنیم.
قبلا برای مقایسه هر دو عددی 
$O(1)$
هزینه می‌دادیم. اکنون با استفاده از روشی که در بالا توضیح داده‌شد دو رشته را مقایسه می‌کنیم.

ثابت کنید در 
$O(Slogn)$
مرحله رشته‌ها مرتب می‌شوند که $S$
برابر با جمع طول رشته‌های داده شده می‌باشد.


%اگر مسئله چند قسمت داشت به این صورت عمل کنید:
%\شروع{زیرمسائل}
%\زیرمسئله[0] % مسئلۀ ۱آ
%\زیرمسئله % مسئلۀ ۱ب
%\زیرمسئله % مسئلۀ ۱پ
%\زیرمسئله % مسئلۀ ۱ت
%\پایان{زیرمسائل}

\solution{
	ثابت می‌کنیم
	ادغام دو تا آرایه از رشته‌ها مثل
	$$s_1, ..., s_k$$
	و
	$$p_1, ..., p_l$$
	با این روش
	$O(\sum_{i=1}^{k} |s_i| + \sum_{i=1}^{l} |p_i|)$
	طول می‌کشد.
	در مرحله اول ادغام
	$s_1$
	و
	$p_1$
	مقایسه می‌شوند که به اندازه
	$O(min(|s_1|, |p_1))$
	طول می‌کشد.
	مستقل از اینکه هر کدام از اینها کوچکتر باشد
	اگر مجموع طول همه رشته‌ها را $S$
	بنامیم چیزی که باقی می‌ماند برابر
	$S - |s_1|$
	یا
	$S - |p_1|$
	است.
	اگر 
	$T(s_1, ..., s_k, p_1, ... , p_l)$
	زمان الگوریتم باشد داریم:
	$$T(s_1, ..., s_k, p_1,..., p_l) \le min(|s_1| , |p_1|) + max(T(s_2, ..., s_k, p_1,...,p_l), T(s_1,...,s_k,p_2,...,p_l)$$
	که در هر صورت نشان می‌دهد:
	$T = O(S)$
	
	پس می‌توان نتیجه گرفت هر لایه از درخت ادغام به اندازه 
	$O(S)$
	طول می‌کشد. می‌دانیم تعداد این لایه‌ها برابر
	$log_2{n}$
	است بنابراین در کل این الگوریتم:
	$O(Slog_2n)$
	طول می‌کشد.
	
}




%%%%%%%%%%%%%%%%%%%%%%%%%%%%%%%%%%%%%%%%%%%%%%%%%%%%%
% برخی دستورات مفید برای تمرین‌ها. %
%%%%%%%%%%%%%%%%%%%%%%%%%%%%%%%%%%%%%%%%%%%%%%%%%%%%%

% فرمول‌های نمونه:
%$f(x) = \Theta(x)$
%$T(x, y) \leq \log(x) + 2^y + \binom{2n}{n}$
% {\tt code\_function}


% لیست موارد بدون شماره‌گذاری:
%\شروع{فقرات}
%    \فقره مورد اول 
%        \شروع{فقرات}
%            \فقره مورد اول داخلی 
%                \شروع{فقرات}
%                    \فقره اولین مورد خیلی داخلی 
%                    \فقره دومین مورد خیلی داخلی 
%                \پایان{فقرات}
%            \فقره مورد دوم داخلی 
%        \پایان{فقرات}
%    \فقره مورد دوم 
%\پایان{فقرات}

% لیست موارد شماره‌گذاری شده:
%\شروع{شمارش}
%    \فقره یک مورد 
%    \فقره یک مورد دیگر 
%    \فقره باز هم یک مورد دیگر
%\پایان{شمارش}

% برای نوشتن یک فرمول چند خطی که یک شماره بخوره. برای حالت‌های دیگر به سایت زیر مراجعه کنید:
% https://www.overleaf.com/learn/latex/Aligning_equations_with_amsmath
%\begin{align} 
%    \begin{split}
%        (x+y)^3 &= (x+y)^2(x+y) \\
%                &= (x^2+2xy+y^2)(x+y) \\
%                &= (x^3+2x^2y+xy^2) + (x^2y+2xy^2+y^3) \\
%                &= x^3+3x^2y+3xy^2+y^3
%    \end{split}                                 
%\end{align}

% نوشتن ماتریس و آرایه دوبعدی
%\begin{align}
%    A = 
%    \begin{bmatrix}
%        A_{11} & A_{21} \\
%        A_{21} & A_{22}
%    \end{bmatrix}
%\end{align}

% اضافه کردن تصویر با شرح. برای اضافه کردن تصویر بدون شرح، فقط دستور درج‌تصویر را استفاده کنید.
%\شروع{شکل}[h!]
%\تنظیم‌ازوسط
%  \درج‌تصویر[پهنا=0.5\پهنای‌متن]{img.jpg}
%\شرح{شرح تصویر}
%\برچسب{تصویر:۱}
%\پایان{شکل}
%\پایان{زیرمسائل}

%کد زیر را ببینید.
%%\lstinputlisting{source_filename.py}
%\begin{latin}
%\begin{lstlisting}
%# Import the modules
%import sys
%import random
%
%ans = True
%
%while ans:
%question = raw_input("Ask the magic 8 ball a question: (press enter to quit) ")
%
%answers = random.randint(1,8)
%
%if question == "":
%sys.exit()
%
%elif answers == 1:
%print "It is certain"
%
%elif answers == 2:
%print "Outlook good"
%
%elif answers == 3:
%print "You may rely on it"
%
%elif answers == 4:
%print "Ask again later"
%
%elif answers == 5:
%print "Concentrate and ask again"
%
%elif answers == 6:
%print "Reply hazy, try again"
%
%elif answers == 7:
%print "My reply is no"
%
%elif answers == 8:
%print "My sources say no"
%\end{lstlisting}
%\end{latin}
